\section{Objetivos}\label{sec:objetivos}

En esta práctica se realizan ejercicios experimentales que sirven de introducción a los algoritmos
cuánticos. A lo largo de ella se adquieren y afianzan conceptos básicos de algoritmia cuántica y se
implementan, con Qiskit, soluciones a los problemas planteados en el enunciado.

Sobre la base de los conocimientos de programación en Qiskit ya consolidados, se construyen los circuitos
de los algoritmos cuánticos y se realizan simulaciones para comparar los resultados experimentales con
las predicciones teóricas.

Los objetivos concretos son:
\begin{itemize}
    \item Estudiar e implementar oráculos cuánticos reversibles a partir de funciones booleanas.
    \item Analizar el efecto de \textbf{retroceso de fase} (\textit{phase kickback}) y su relación con los oráculos de fase.
    \item Implementar y analizar los algoritmos de consulta \textbf{Deutsch}, \\
    \textbf{Deutsch-Jozsa} y \textbf{Bernstein-Vazirani}.
    \item Implementar y analizar el \textbf{algoritmo de Grover}, incluido el número óptimo de iteraciones.
    \item Implementar la \textbf{estimación cuántica de fase (QPE)} de un operador $U$ sobre un autovector $\ket{\psi}$.
    \item Contrastar en cada caso los resultados simulados con los teóricos.
\end{itemize}